\chapter{การออกแบบและประกอบชิ้นส่วนฮาร์ดแวร์}
\label{chap:ch02}

\section*{แผนบริหารการสอนประจำบทที่ 2}
\addcontentsline{toc}{section}{แผนบริหารการสอนประจำบทที่ 2}

\begin{planbox}{หัวข้อเนื้อหาประจำบท}
\begin{enumerate}[leftmargin=1.5em, itemsep=1pt]
    \item สถาปัตยกรรมฮาร์ดแวร์และคุณสมบัติของไมโครคอนโทรลเลอร์ ATSAMD51
    \item รายการอุปกรณ์และโมดูลตรวจวัดแสง TCS34725 ร่วมกับหลอดไฟ WS2812B
    \item การออกแบบทางเดินแสงของช่องใส่คิวเวตต์เพื่อลดการสะท้อนและการกระเจิง
    \item ขั้นตอนการประกอบกลไกและการต่อขาสัญญาณดิจิทัล
\end{enumerate}
\end{planbox}

\begin{planbox}{วัตถุประสงค์เชิงพฤติกรรม}
\begin{enumerate}[leftmargin=1.5em, itemsep=1pt]
    \item ระบุหน้าที่และตำแหน่งของขาสัญญาณควบคุมในระบบสเปกโทรโฟโตมิเตอร์ได้อย่างถูกต้อง
    \item อธิบายหลักการจัดวางแนวแกนแสงระหว่างแหล่งกำเนิด หลอดคิวเวตต์ และเซนเซอร์ได้อย่างสมบูรณ์
    \item ปฏิบัติการประกอบและทดสอบการเชื่อมต่อวงจรฮาร์ดแวร์ได้อย่างถูกต้องปลอดภัย
\end{enumerate}
\end{planbox}

\begin{planbox}{กิจกรรมการเรียนการสอน}
\begin{enumerate}[leftmargin=1.5em, itemsep=1pt]
    \item การบรรยายและสาธิตการเชื่อมต่อสายสัญญาณผ่านพอร์ต Grove
    \item ปฏิบัติการประกอบชิ้นส่วนฮาร์ดแวร์ลงในกล่องบรรจุ 3D Printed Enclosure
    \item การตรวจสอบความถูกต้องของระบบไฟฟ้าและสัญญาณ I2C ก่อนการจ่ายไฟ
\end{enumerate}
\end{planbox}

\begin{planbox}{สื่อการเรียนการสอน}
\begin{enumerate}[leftmargin=1.5em, itemsep=1pt]
    \item บอร์ด Seeed Studio Wio Terminal พร้อมสายเชื่อมต่อ Grove
    \item โมดูลเซนเซอร์ Adafruit TCS34725 และโมดูลไฟ Grove NeoPixel WS2812B
    \item ชุดโครงสร้างกล่องบรรจุ 3 มิติและแท่นยึดคิวเวตต์ขนาด 10 มิลลิเมตร
\end{enumerate}
\end{planbox}

\begin{planbox}{การวัดและประเมินผล}
\begin{enumerate}[leftmargin=1.5em, itemsep=1pt]
    \item การประเมินทักษะการประกอบวงจรและการจัดสายสัญญาณอย่างเป็นระเบียบ
    \item การทดสอบการตรวจพบอุปกรณ์ I2C ผ่านโปรแกรมตรวจสอบที่อยู่บัส
    \item การตรวจความถูกต้องของการตอบคำถามทบทวนท้ายบทที่ 2
\end{enumerate}
\end{planbox}

\clearpage

\section{รายการอุปกรณ์และโมดูลฮาร์ดแวร์}
อุปกรณ์ที่ใช้ในการประกอบและสร้างเครื่องสเปกโทรโฟโตมิเตอร์ประกอบด้วยรายการดังแสดงในตารางที่ \ref{tab:components}

\begin{table}[H]
    \caption{รายการอุปกรณ์และโมดูลฮาร์ดแวร์}
    \label{tab:components}
    \begin{tabularx}{\textwidth}{p{3.2cm}p{4.2cm}X}
        \toprule
        \textbf{อุปกรณ์} & \textbf{รุ่นหรือรหัสสินค้า} & \textbf{บทบาทหน้าที่ในระบบ} \\
        \midrule
        บอร์ดประมวลผล & Seeed Studio Wio Terminal & ไมโครคอนโทรลเลอร์ ATSAMD51 พร้อมจอ LCD 2.4 นิ้ว \\
        เซนเซอร์ตรวจวัดแสง & Adafruit TCS34725 & เซนเซอร์ตรวจวัดแสงสีและค่าความสว่างแบบ I2C \\
        แหล่งกำเนิดแสงสี & Grove RGB LED (WS2812B) & ไดโอดเปล่งแสงควบคุมสีแบบดิจิทัลความยาวคลื่นแม่นยำ \\
        ช่องใส่หลอดทดลอง & Cuvette Holder 10 mm & แท่นยึดหลอดคิวเวตต์มาตรฐานทางเดินแสง 10 มิลลิเมตร \\
        สายเชื่อมต่อข้อมูล & Grove 4-Pin Cable & สายส่งสัญญาณดิจิทัลและสายจ่ายพลังงานไฟฟ้า \\
        การ์ดหน่วยความจำ & microSD Card Class 10 & อุปกรณ์บันทึกผลการวัดลงไฟล์ข้อมูล 4 รูปแบบ (NPK.csv, SPECTRUM.csv, CALIB\_LOG.csv, LIQUID\_LOG.csv) \\
        โครงสร้างตัวเครื่อง & 3D Printed Enclosure & กล่องบรรจุอุปกรณ์ภายนอกออกแบบด้วยเส้นใย PLA+ \\
        \bottomrule
    \end{tabularx}
\end{table}

\section{แผนภาพการต่อวงจรและขาสัญญาณ}
การเชื่อมต่อสัญญาณระหว่างบอร์ด Wio Terminal กับโมดูลเซนเซอร์และแหล่งกำเนิดแสง แสดงรายละเอียดตามตารางที่ \ref{tab:pinout}

\begin{table}[H]
    \caption{การกำหนดขาสัญญาณและการเชื่อมต่ออุปกรณ์}
    \label{tab:pinout}
    \begin{tabularx}{\textwidth}{p{3.0cm}p{3.0cm}p{3.0cm}X}
        \toprule
        \textbf{อุปกรณ์ภายนอก} & \textbf{พอร์ต Wio Terminal} & \textbf{ขาสัญญาณไมโคร} & \textbf{คำอธิบายการเชื่อมต่อ} \\
        \midrule
        TCS34725 (SDA) & Grove พอร์ตขวา & ขา SDA (PA17) & ส่งสัญญาณข้อมูล I2C บัสความเร็ว 100 kHz \\
        TCS34725 (SCL) & Grove พอร์ตขวา & ขา SCL (PA16) & ส่งสัญญาณนาฬิกา I2C บัสความเร็ว 100 kHz \\
        TCS34725 (VCC) & Grove พอร์ตขวา & ขา 3V3 หรือ 5V & แหล่งจ่ายพลังงานไฟฟ้ากระแสตรง 3.3 โวลต์ \\
        TCS34725 (GND) & Grove พอร์ตขวา & ขา GND & กราวด์ร่วมของระบบ \\
        \midrule
        RGB LED (DIN)  & Grove พอร์ตซ้าย & ขา D0 (PB08) & ขาสัญญาณข้อมูลควบคุมหลอด NeoPixel WS2812B \\
        RGB LED (VCC)  & Grove พอร์ตซ้าย & ขา 5V        & จ่ายแรงดันไฟฟ้าเลี้ยงหลอดไฟ 5 โวลต์ \\
        RGB LED (GND)  & Grove พอร์ตซ้าย & ขา GND       & กราวด์ร่วมของหลอดไฟ \\
        \midrule
        ปุ่มกด A       & ด้านบนตัวเครื่อง & WIO\_KEY\_A (28) & ปุ่มด้านขวาสุด ควบคุมการเปิดปิดแสงสีแดง \\
        ปุ่มกด B       & ด้านบนตัวเครื่อง & WIO\_KEY\_B (29) & ปุ่มตรงกลาง ควบคุมการเปิดปิดแสงสีเขียว \\
        ปุ่มกด C       & ด้านบนตัวเครื่อง & WIO\_KEY\_C (30) & ปุ่มด้านซ้ายสุด ควบคุมการเปิดปิดแสงสีน้ำเงิน \\
        จอยสติ๊กกลาง   & ด้านหน้าตัวเครื่อง & WIO\_5S\_PRESS (35) & กดลงตรงกลาง ควบคุมการเปิดปิดแสงสีขาว \\
        \bottomrule
    \end{tabularx}
\end{table}

\section{การออกแบบทางเดินแสงและกล่องบรรจุ 3 มิติ}
กล่องบรรจุอุปกรณ์ภายนอก (Enclosure) และแท่นยึดคิวเวตต์ได้รับการขึ้นรูปด้วยเครื่องพิมพ์สามมิติ (3D Printing) โดยใช้วัสดุเส้นใย PLA+ สีดำด้าน (Matte Black) เพื่อลดการสะท้อนแสงภายในห้องวัด (Internal Reflection Reduction) โดยมีคุณลักษณะทางทัศนศาสตร์ที่สำคัญ
\begin{itemize}[leftmargin=1.5em, itemsep=2pt]
    \item \textbf{แนวแกนแสงตรง (Coaxial Optical Path):} กึ่งกลางหลอดไฟ WS2812B ช่องคิวเวตต์ และเซนเซอร์ TCS34725 อยู่ในระดับความสูงเดียวกันที่ 15.0 มิลลิเมตรจากฐาน
    \item \textbf{ช่องจำกัดลำแสง (Collimating Aperture):} ขนาดเส้นผ่านศูนย์กลาง 3.0 มิลลิเมตร ติดตั้งระหว่างหลอด LED กับผิวคิวเวตต์ เพื่อให้ได้ลำแสงขนานที่ตั้งฉากกับผนังคิวเวตต์
    \item \textbf{ฝาครอบกันแสงภายนอก (Light-Tight Lid):} มีขอบประกบแบบเขี้ยวซ้อน (Baffle Joint) ป้องกันแสงรอบข้าง (Ambient Light Stray) เล็ดลอดเข้าสู่เซนเซอร์ได้อย่างสมบูรณ์
\end{itemize}

\section{ขั้นตอนการประกอบชิ้นส่วนกลไก}
การประกอบอุปกรณ์ลงในกล่องบรรจุ 3D Printed Enclosure มีขั้นตอนดังต่อไปนี้

\begin{enumerate}[leftmargin=1.5em, itemsep=3pt]
    \item นำแท่นยึดหลอดคิวเวตต์ (Cuvette Holder) ติดตั้งเข้ากับฐานล่างของกล่อง โดยหันด้านช่องส่องแสงให้ตรงกับตำแหน่งติดตั้งโมดูลหลอดไฟและเซนเซอร์
    \item ติดตั้งโมดูลหลอดไฟ Grove RGB LED เข้ากับด้านข้างช่องคิวเวตต์ โดยให้ผิวหน้าของหลอด LED ขนานกับผนังคิวเวตต์และกั้นแสงรบกวนจากภายนอก
    \item ติดตั้งโมดูลเซนเซอร์ TCS34725 ในตำแหน่งตรงข้ามกับหลอด LED ในแนวแกนแสงเดียวกัน เพื่อให้รับลำแสงที่ทะลุผ่านหลอดคิวเวตต์ได้อย่างสมบูรณ์
    \item เสียบสาย Grove 4-Pin จากหลอด LED เข้าสู่พอร์ต Grove ด้านซ้าย (D0) ของบอร์ด Wio Terminal
    \item เสียบสาย Grove 4-Pin จากเซนเซอร์ TCS34725 เข้าสู่พอร์ต Grove ด้านขวา (I2C) ของบอร์ด Wio Terminal
    \item ประกอบบอร์ด Wio Terminal เข้ากับฝาครอบด้านหน้า ขันสกรูยึดให้แน่นหนา และตรวจเช็คการขยับของปุ่มกด A, B, C และจอยสติ๊กด้านหน้าให้เคลื่อนไหวได้อย่างอิสระ
\end{enumerate}

\section*{บทสรุปประจำบทที่ 2}
\addcontentsline{toc}{section}{บทสรุปประจำบทที่ 2}
การออกแบบโครงสร้างฮาร์ดแวร์ของเครื่องสเปกโทรโฟโตมิเตอร์คำนึงถึงความแม่นยำทางทัศนศาสตร์และความสะดวกในการพกพา การใช้พอร์ต Grove อำนวยความสะดวกในการเชื่อมต่อเซนเซอร์ I2C และหลอดดิจิทัล NeoPixel อย่างแน่นหนา ในขณะที่ห้องวัดแสงที่พิมพ์ด้วยเส้นใยสีดำด้านพร้อมช่องจำกัดลำแสงช่วยลดสัญญาณรบกวนและแสงกระเจิง ส่งผลให้สัญญาณการวัดมีความเสถียรและทำซ้ำได้สูง

\section*{คำถามทบทวนท้ายบทที่ 2}
\addcontentsline{toc}{section}{คำถามทบทวนท้ายบทที่ 2}
\begin{enumerate}[leftmargin=1.5em, itemsep=3pt]
    \item เหตุใดการออกแบบแท่นยึดคิวเวตต์จึงต้องใช้วัสดุสีดำด้าน และหากใช้วัสดุสีขาวสะท้อนแสงจะส่งผลกระทบต่อค่าการดูดกลืนแสงอย่างไร
    \item สัญญาณนาฬิกา I2C ที่ความเร็ว 100 kHz เหมาะสมกับการอ่านค่าจากเซนเซอร์ TCS34725 หรือไม่ เพราะเหตุใด
    \item จงอธิบายบทบาทของช่องจำกัดลำแสง (Collimating Aperture) ที่อยู่ระหว่างหลอดไฟ LED กับคิวเวตต์
    \item หากเซนเซอร์ TCS34725 ไม่อยู่ในแนวแกนเดียวกับหลอด LED จะเกิดความคลาดเคลื่อนเชิงแสงในลักษณะใด
\end{enumerate}
